\section{Four boxes and beyond}\label{sec:higher}

The tables of KN stop at three boxes. In this section we test all $\rho$ with $4\le|\rho|\le9$, in ranges depending on the knot, for torus knots, for which the refined invariants are defined.
We also describe a weak test for $\mathbf{4_1}$ with four boxes in which the unknown $\rCS_{\yng(3,1)}$ is constrained, and analyse the structure of the KN candidate for $\yng(2,1)$.
Everything in this section is at framing zero.

\subsection{Computing \texorpdfstring{$\hf_\rho$}{f-hat} for any number of boxes}
The matrix $M(t)$ in the relation \eqref{eq:f-to-hatf} is written explicitly in KN for $|\mu|\le3$. For general size we use the following diagonalisation.
% [M1]
\begin{lemma}\label{lem:M-diag}
For $|\mu|=|\rho|=n$, the matrix $M_{\mu\rho}(t)=\sum_\sigma C_{\mu\sigma\rho}B_\sigma(t)$ is the matrix of the Kronecker product with $B(t)=\sum_\sigma B_\sigma(t)s_\sigma$,
where $C_{\mu\sigma\rho}$ are the Kronecker coefficients. It is diagonal in the basis of characters of the symmetric group:
\[
\sum_\mu M_{\mu\rho}(t)\,\chi_\mu(\kappa)=B(\kappa)\,\chi_\rho(\kappa),\qquad
B(\kappa)=t^{-\frac{n-1}2}\,\frac{\prod_i(1-t^{\kappa_i})}{1-t},
\]
where $\kappa$ runs over conjugacy classes (partitions of $n$).
\end{lemma}
\begin{proof}
By $C_{\mu\sigma\rho}=\sum_\nu z_\nu^{-1}\chi_\mu(\nu)\chi_\sigma(\nu)\chi_\rho(\nu)$ and the orthogonality of characters, the left-hand side equals $\chi_\rho(\kappa)\sum_\sigma B_\sigma(t)\chi_\sigma(\kappa)$.
$B_\sigma$ is non-zero only for hooks $\sigma=(n-d,1^d)$, where it equals $(-t)^d t^{-(n-1)/2}$,
and $\sum_{d}(-u)^d s_{(n-d,1^d)}=\sum_\kappa z_\kappa^{-1}p_\kappa\prod_i(1-u^{\kappa_i})/(1-u)$.
\end{proof}
Hence $\hf_\rho$ is obtained by dividing the coefficients of $f^q$ in the power-sum basis by $B(\kappa)$ and returning to the Schur basis with the characters.
The plethystic logarithm is computed at each evaluation point by a degree-by-degree recursion in the power-sum basis.
The invariants $\rCS_\lambda$ are the elliptic Hall algebra values multiplied by the closed normalisation $M_\lambda=M_\square^{|\lambda|}T_\lambda^{mn}$ of Computational Result~\ref{cr:M}.
This procedure reproduces in all coefficients the three-box $\hf_\rho$ computed with the explicit formulae (for $T_{2,3}$ and $T_{3,4}$, all six $\rho$),
and the consistency relation $f^{\bar t}=M(q^{-1})\hf$ holds at all evaluation points.

\subsection{Torus knots}
% [R18]
\begin{compresult}\label{cr:higher}
For the eight torus knots and the numbers of boxes in Table~\ref{tab:higher}, $\hf_\rho$ is positive for all $\rho$ at framing zero.
The number of pairs $(K,\rho)$ with $|\rho|\ge4$ tested is $340$.
\end{compresult}
\begin{table}[h]
\centering
\caption{Tests with four or more boxes at framing zero. ``Cases'' is the total number of $\rho$ tested; ``terms'' is the sum over $\rho$ of the numbers of terms of the spin form at the largest number of boxes.
There are no negative terms in the spin form or in the genus form.}\label{tab:higher}
\begin{tabular}{lcrrr}
\toprule
Knot & $|\rho|$ & cases & terms (largest $|\rho|$) & time for the largest $|\rho|$ (s) \\
\midrule
$T_{2,3}$ & $4$--$9$ & 90 & 1786668 & 428 \\
$T_{2,5}$ & $4$--$8$ & 60 & 2880410 & 299 \\
$T_{2,7}$ & $4$--$8$ & 60 & 6556642 & 1260 \\
$T_{2,9}$ & $4$--$8$ & 60 & 11740842 & 2036 \\
$T_{3,4}$ & $4$--$6$ & 23 & 1358553 & 280 \\
$T_{3,5}$ & $4$--$6$ & 23 & 2336943 & 322 \\
$T_{3,7}$ & $4$--$5$ & 12 & 1307758 & 64 \\
$T_{4,5}$ & $4$--$5$ & 12 & 1848265 & 474 \\
\bottomrule
\end{tabular}
\end{table}
The ranges were chosen as follows.
At four boxes the non-hook colour $\yng(2,2)$ and the Adams images $(a,q,t)\mapsto(a^2,q^2,t^2)$ of $\yng(2)$ and $\yng(1,1)$ enter $\hf_\rho$ for the first time.
At six boxes the Adams operations with $d=2,3,6$ have appeared for all colours with at most three boxes.
From seven boxes on no new type of term appears and only the size grows.
The knots were chosen by the number of strands, i.e.\ the width in the $a$-degree: $T_{2,2k+1}$ are thin, $T_{3,n}$ are thick of width $2$, and $T_{4,5}$ has width $3$.
\begin{remark}
The overall sign was $-1$ for $|\rho|\equiv2\pmod4$ and $+1$ otherwise, uniformly for all $\rho$ with the same number of boxes.
For odd $|\rho|$ the half-integer shift of the $a$-degree is the periodic continuation of the convention of KN for $|\rho|=1,3$; it shifts $\beta$ uniformly and does not affect positivity.
The coefficients of the genus form grow rapidly with the number of boxes: up to about $10^{30}$ for $T_{2,3}$ with nine boxes, and already about $10^{21}$ for $T_{4,5}$ with three boxes.
\end{remark}

\subsection{The figure-eight knot with four boxes: a weak test}\label{sec:fig8-four}
For $\mathbf{4_1}$ with four boxes the rectangular inputs $\yng(4)$, $\yng(2,2)$, $\yng(1,1,1,1)$ come from the refined formulae of KNTZ \cite{Kameyama19};
they agree in all coefficients with the values of KN for $\yng(1)$ to $\yng(1,1,1)$ and with the GNSSS formulae for the symmetric colours $[r]$, $r\le4$.
For $\yng(2,1)$ we use the KN candidate.
The refined value for $\yng(3,1)$, and for its transpose $\yng(2,1,1)$, is not in the literature.
We therefore treat the coefficients of $\rCS_{\yng(3,1)}$ as unknowns subject to the following linear conditions:
\begin{itemize}
\item at $t=q$ it equals the differential expansion of Morozov \cite{Morozov16} for $\yng(3,1)$ (for $\mathbf{4_1}$, a sum of $Z$-factors);
\item transposition, $\rCS_{\yng(2,1,1)}(a,q,t)=\rCS_{\yng(3,1)}(a,t^{-1},q^{-1})$, and amphichirality;
\item exponential growth: it equals $\rCS_{\yng(1,1)}\rCS_{\yng(1)}^2$ at $q=1$ and $\rCS_{\yng(3)}\rCS_{\yng(1)}$ at $t=1$;
\item at $a=t^2$ it equals $\rCS_{\yng(2)}\rCS_{\yng(1,1)}$ (reduction to $U(2)$; a hypothesis);
\item all $\hf_\rho$ with $|\rho|=4$ are Laurent polynomials, the genus expansion closes, and they are positive.
\end{itemize}
Since $\hf_\rho$ is linear in $\rCS_{\yng(3,1)}$, with a denominator depending only on $q,t$, these are conditions of a linear programme.
% [R19]
\begin{compresult}\label{cr:fig8-four}
With the unknowns restricted to $2367$ monomials determined from the values at the specialisations, the following hold.
\begin{itemize}
\item With overall sign $-1$ the problem is infeasible, already with positivity of the spin form alone.
\item With overall sign $+1$ there exists a rational solution for which both forms are non-negative.
\item No integral solution was found by mixed integer programming within $50$ minutes; its existence is undecided.
\end{itemize}
\end{compresult}
This is not a test of the conjecture; it only shows that, over the rationals, an $\rCS_{\yng(3,1)}$ compatible with the conjecture exists at four boxes.
The conclusion depends on the KN candidate for $\yng(2,1)$, on the $U(2)$ hypothesis, and on the range of unknowns.
The formulation was checked on $T_{2,3}$, for which $\rCS_{\yng(3,1)}$ is known from the elliptic Hall algebra:
the true value satisfies all conditions and lies in the range of unknowns, and the $\hf_\rho$ assembled from it agree with the direct computation.

\subsection{The structure of the KN candidate for \texorpdfstring{$\protect\yng(2,1)$}{[2,1]}}\label{sec:fig8-21}
KN state that their candidate for $\mathbf{4_1}$ coloured by $\yng(2,1)$ was found by imposing mirror symmetry, exponential growth and amphichirality \cite{Kameyama17}.
It can be written in terms of the $Z$-factor expansion of Anokhina et al.\ \cite{Anokhina12}.
Let $(A,Q,T)=(a,q,t)^{1/2}$, $\{z\}=z-z^{-1}$ and $\fZ_{I|J}^{(s|\sigma)}=\{AQ^{I+s}T^{-\sigma}\}\{AQ^sT^{-J-\sigma}\}$, with $\fZ_{I|J}=\fZ_{I|J}^{(0|0)}$.
% [R20]
\begin{compresult}\label{cr:fig8-21}
The KN candidate equals
\begin{multline*}
1+\bigl(\fZ^{(-1|-1)}_{3|3}+\fZ^{(1|0)}_{1|1}+\fZ^{(0|1)}_{1|1}\bigr)
+\bigl(\fZ^{(0|-1)}_{3|3}\fZ^{(1|0)}_{1|1}+\fZ^{(-1|0)}_{3|3}\fZ^{(0|1)}_{1|1}+\alpha\,\fZ^{(1|0)}_{1|1}\fZ^{(0|1)}_{1|1}\bigr)\\
+\fZ_{3|3}\fZ^{(1|0)}_{1|1}\fZ^{(0|1)}_{1|1}
\end{multline*}
with $\alpha=1-\{Q\}\{T\}$.
Anokhina et al.\ take the same expression with $\alpha=1-(Q-T^{-1})(T-Q^{-1})$ and obtain a superpolynomial with all coefficients positive.
In bold variables the KN candidate has $69$ terms of which $10$ are negative, and its value at $\ba=\bq=\bt=1$ is $125=5^3$, the cube of the value for $\yng(1)$; for the expression of Anokhina et al.\ the value is $189$.
\end{compresult}
At $t=q$ both choices of $\alpha$ reduce to $1-(Q-Q^{-1})^2$.
Thus the KN candidate corresponds to imposing exponential growth instead of positivity of the superpolynomial.

We also examined the condition of Heegaard--Floer type proposed by Anokhina et al.\ for hook colours $\lambda=[r,1^{s-1}]$:
\begin{equation}\label{eq:HF}
\rCS_\lambda\bigl(A=T/Q\,\big|\,Q,T\bigr)=\rCS_{\yng(1)}\bigl(A=T/Q\,\big|\,Q^rT^{s-1},\,T^sQ^{r-1}\bigr).
\end{equation}
% [R21]
\begin{compresult}\label{cr:HF}
\eqref{eq:HF} holds for $\mathbf{4_1}$ with $\yng(2,1)$, both for the KN candidate and for the expression of Anokhina et al.
For the elliptic Hall algebra values of $T_{2,3}$, $T_{2,5}$, $T_{3,4}$ with $\yng(2,1)$ and $\yng(3,1)$ it holds up to the monomial $(T/Q)^{\delta(|\lambda|-1)}$.
\end{compresult}
Condition \eqref{eq:HF} is a linear condition that was not included in the linear programme above.

Finally, we tried to extend the method of KN to $\yng(3,1)$ by decomposing the value of Morozov at $t=q$ into the form of Anokhina et al.:
one product of shifted $Z$-factors for each subset of boxes, with corrections of the coefficients only for disconnected subsets.
For $\yng(2,1)$ this search reproduces the expansion of Anokhina et al.\ as the unique solution.
For $\yng(3,1)$ no decomposition of this form was found (the integer programme with shifts $|k|\le2$ is infeasible, and there is no solution when the shifts are sums of contributions of pairs of boxes).
The colour $\yng(3,1)$ seems to require structure beyond products over subsets.
